\documentclass[11pt]{article}
%autocompile publish
\usepackage{math-hse,array}
\title{Лекция 12. Элементы теории вероятностей, часть 1}

\date{19.11.2025}
\begin{document}
\section{Основные понятия}
\paragraph{Множество элементарных исходов} Пусть должно произойти испытание, исход которого неизвестен, но берётся из какого-то определённого, известного нам множества. Классические примеры ситуаций: я брошу игральную кость, что на ней выпадет? Посмотрим завтра в газеты, какая будет цена на нефть? Выйдем на улицу 15 января 2050 года, какая будет температура воздуха? Мы используем здесь слово <<испытание>> или <<опыт>>, но совершенно не обязательно речь идёт о каких-то научных экспериментах~--- мы испытываем окружающий мир с разных сторон. Множество всех возможных исходов случайного эксперимента, неразложимых на более простые и попарно несовместных, будем называть \emph{множеством элементарных исходов}.

\begin{examples}\hspace{1em}\\
\begin{enumerate}
	\item Если мы кидаем игральную кость, то пространством элементарных исходов можно считать числа от 1 до 6.

	\item Если мы обсуждаем температуру 15 января 2050 года, то пространством элементарных исходов можно считать числа от --100 до +100.

	\item Исходом случайного испытания необязательно должно быть число. Что выпадет при бросании монетки: орёл или решка? На кого из ребят падёт жребий водить в игре: на Саню, Петю, Митю или Лену?
\end{enumerate}
\end{examples}

Здесь мы вплотную встречаемся с элементами математического моделирования. Разве мы перечислили все исходы при бросании кубика? Конечно, нет. Ведь кубик мог лететь высоко или низко, упасть далеко или близко, задеть чей-то лоб, закатиться под стол и~т.\,д. и~т.\,п. Нас же интересует только, какое число выпало на верхней грани кубика.

\emph{События и благоприятные исходы} Множество всех элементарных исходов обычно обозначается буквой $\Omega$. \emph{Событием} мы будем называть любое подмножество множества элементарных исходов. Если элементарный исход принадлежит какому-то событию, мы будем говорить, что элементарный исход \emph{благоприятствует} этому событию, элементарный исход \emph{благоприятный} для этого события.

\begin{examples}\hspace{1em}\\
\begin{enumerate}
	\item Случайное испытание: что выпадет на игральной кости? Случайное событие $A$ ~--- на кубике выпало чётное число ~---    описывается таким подмножеством множества $\Omega: \left\{2,4,6\right\}$. Исход бросания кубика <<2>> благоприятствует событию $A$, это благоприятный исход для выпадения чётного числа.
	
	\item Случайное испытание: какая будет погода 15 января 2050 года? Случайное событие $B$~--- на улице 15 января 2050 года холодно~--- описывается таким подмножеством множества $\Omega: \left\{\left[-100;-5\right]\right\}$. Исход измерения температуры $-20$~благоприятствует событию $B$, а $+10$~--- не благоприятствует.
	
	\item Случайное испытание: троекратное бросание монетки. Случайное событие $C$~--- выпадение хотя бы двух орлов. Элементарный исход <<орёл, решка, орёл>> благоприятствует событию $C$, а элементарный исход <<решка, решка, орел>> не благоприятствует событию $C$.
\end{enumerate}
\end{examples}

Мы видим, что обычное, <<бытовое>> понятие события описывается на языке множеств"--~подмножеств.

\paragraph{Несовместные события} Если два события не могут произойти одновременно, мы будем называть их \emph{несовместными}. На языке подмножеств это означает, что не найдется ни одного элементарного исхода, принадлежащего обоим событиям.

\begin{examples}\hspace{1em}\\
\begin{enumerate}
	\item При бросании игральной кости события $A$~--- выпадение чётного числа и $B$~--- выпадение числа, меньшего трёх ~--- совместны, т.\,к. элементарный исход <<выпадение двойки>> \ благоприятен для обоих событий. Заметим, что не все элементарные исходы, благоприятные $A$, благоприятны и $B$: скажем, элементарный исход <<выпадение четвёрки>> благоприятен событию $A$, но не благоприятен $B$.

	\item Случайное испытание: троекратное бросание монетки. Случайное событие $A$ ~--- выпадение хотя бы двух орлов и $B$ ~--- выпадение хотя бы одной решки. Эти события совместны: элементарный исход <<решка, орёл, орёл>> благоприятствует обоим событиям. Предположим, событие $C$ заключается в том, что количество выпавших решек больше, чем количество выпавших орлов. Тогда события $A$ и событие $C$ несовместны.
\end{enumerate}
\end{examples}

\paragraph{Пересечение событий} Пусть $A$ и $B$ некоторые события. Событие $C$, состоящее в том, что выполняются события $A$ и $B$, мы будем называть \emph{пересечением} событий $A$ и $B$ и обозначать $AB$ или $A\cap B$.

\begin{examples}
	\item При бросании игральной кости события $A$ ~--- выпадение чётного числа и $B$ ~--- выпадение числа, меньшего трёх. Событие $A\cap B$ состоит в выпадении чётного числа меньшего трёх, т.\,е. в выпадении двойки.

	\item Случайное испытание: троекратное бросание монетки. Случайное событие $A$~--- выпадение хотя бы двух орлов и $B$~--- выпадение хотя бы одной решки. Событие $A\cap B$ состоит в выпадении двух орлов и одной решки, ему благоприятны следующие элементарные исходы: <<орёл, орёл, решка>>, <<орёл, решка, орёл>> и <<решка, орёл, орёл>>.
\end{examples}

\paragraph{Объединение событий} Пусть $A$ и $B$ некоторые события. Событие $C$, состоящее в том, что выполняются события $A$ или $B$ (можно оба сразу), мы будем называть \emph{объединением} событий $A$ и $B$ и обозначать $A+B$ или $A\cup B$.

Рассмотрим такую ситуацию, в которой все исходы <<равноправны>>, то есть встречаются одинаково часто, если повторять эксперимент много раз. Скажем, элементарные исходы <<1>>, <<2>>, <<3>>, <<4>>, <<5>>, <<6>> при бросании игральной кости естественно считать равноправными (если, конечно, кость не кривая), а вот в случайном испытании <<Встречу ли я динозавра, выходя сегодня из дверей университета?>> исходы <<да>> и <<нет>> равноправными не являются.

\paragraph{Вероятность события} \emph{Вероятностью события A} в случае равноправности элементарных исходов, мы будем называть отношение количества элементарных исходов, благоприятных событию $A$ к общему количеству элементарных исходов:
$$
	P(A)=\frac{N(A)}{N(\Omega)}
$$

\begin{examples}
	\item Найдём вероятность выпадения чётного числа на игральной кости: $P(A)=N(A)/N(\Omega)=3/6=1/2$, т.\,к. ~на игральной кости всего три чётных числа: ~2, ~4, ~6.

	\item Найдём вероятность выпадения хотя бы одной решки при троекратном бросании монетки.

Вот все возможные исходы, их восемь:

   <<орёл, орёл, орёл>>

   <<орёл, орёл, решка>>

   <<орёл, решка, орёл>>

   <<орёл, решка, решка>>

   <<решка, орёл, орёл>>

   <<решка, орёл, решка>>

   <<решка, решка, орёл>>

   <<решка, решка, решка>>

Событию <<выпала хотя бы одна решка>> благоприятны семь исходов, все кроме первого. Итак, искомая вероятность:
$P(A)=N(A)/N(\Omega)=7/8$.
\end{examples}

Заметим, что \emph{сумма вероятностей всех элементарных исходов случайного эксперимента равна единице}:
$$
	N(\Omega)\times\frac{1}{N(\Omega)}=1.
$$

\begin{example}
Случайный эксперимент: кидаем игральную кость. Имеем шесть возможных элементарных исходов, вероятность каждого $1/6$, сумма вероятностей равна 1.
\end{example}

\paragraph{Дополнительные события} Мы будем говорить, что события $A$ и $B$ являются \emph{дополнительными}, если $A$ и $B$ вместе составляют все пространство элементарных исходов $\Omega$ и являются несовместными. Дополнение к событию $A$ обозначают $B=\overline{A}$.

\begin{example}
	События $A$ ~--- на игральной кости выпало чётное число  и $B$ ~--- на игральной кости выпало нечётное число ~--- являются дополнительными: действительно, каждое выпавшее на игральной кости число будет обязательно или чётным, или нечётным.
\end{example}

\emph{Сумма вероятностей дополнительных событий равна единице}.
\begin{examples}
   \item Пусть события $A$ ~--- на игральной кости выпало чётное число и $B$ ~--- на игральной кости выпало нечётное число. Тогда $P(A)=3/6=1/2$, т.\,к. ~событию $A$ благоприятны три исхода: выпадение <<2>>, <<4>> или <<6>>; $P(B)=3/6=1/2$, т.\,к. ~событию $B$ благоприятны три исхода: выпадение <<1>>, <<3>> или <<5>>. Сумма вероятностей $P(A)+P(B)=1/2+1/2=1$.\\

\end{examples}

\end{document}
